\documentclass{amsart}

\input{../preamble_general.tex}

%\author{Jalen Chrysos}
\title{Practice Midterm 1 (with Solutions)}

% topics:
% induction
% counting 
% induced subgraphs
% graph isomorphism, "triangle free"
% graph score and Havel-Hakimi
% K_n, K_{m,n} and C_n
% connected components
% graph distance
% 2-connected graphs, ear decomposition


\begin{document}
	
	\maketitle 
	
	\textit{In all questions, please justify your answers as appropriate. Any legitimate solution will be accepted, unless the question specifically calls for a certain method.}\\
	
	\noindent \textbf{Question 1}: 
	\begin{enumerate}[(a)]
		\item Show that there is a graph $G$ whose degree sequence (score) is $(3,3,3,3,3,5,5,5)$. \\
		
		One way to quickly do this is to see that $K_{3,5}$ has this degree sequence. Otherwise, there is the Havel-Hakimi algorithm:
		\begin{align*}
			(3,3,3,3,3,5,5,5) &\to (3,3,2,2,2,4,4)\\
			&\to (2,2,2,3,3,4,4) \\
			&\to (2,2,1,2,2,3) \\
			&\to (1,2,2,2,2,3) \\
			&\to (1,2,1,1,1)\\
			&\to (1,1,1,1,2)\\
			&\to (1,1,0,0)\\
			&\to (0,0,1,1)\\
			&\to (0,0,0).
		\end{align*}
		
		\item Let $G$ be a graph as in part (a). How many edges must $G$ contain? \\
		
		The sum of the degrees of a graph is equal to twice the number of edges (handshake lemma). Thus, the number of edges is 
		$$
		\frac{1}{2}(3+3+3+3+3+5+5+5) = 15.
		$$
		\vspace{50px}
		
		\item Is it possible that such a graph $G$ contains $C_6$ as an induced subgraph? \\
		
		No. Suppose $C_6$ is an induced subgraph of $G$ and the other two vertices are labeled $a,b$. Then each vertex in the $C_6$ has degree at most 4: it cannot connect to any of the three vertices of the $C_6$ that are not adjacent to it, because the $C_6$ is \textit{induced}. However, the graph $G$ is supposed to have three vertices of degree 5. There cannot be three vertices of degree 5 and six vertices of degree at most 4, since there are 8 vertices total.
	\end{enumerate}
	
	\newpage 
	
	\noindent \textbf{Question 2}: In this question it is okay to leave your answer unsimplified.
	\begin{enumerate}[(a)]
		\item How many subgraphs of $K_{10}$ are isomorphic to $K_4$? \\
		
		Each $K_4$ subgraph in $K_{10}$ is the induced subgraph of a selection of $4$ vertices, of which there are ${10 \choose 4}$.
		\vspace{70px}
		
		\item How many subgraphs of $K_{10}$ are isomorphic to $K_{3,3}$?\\
		
		\noindent There are ${10\choose 6}$ ways to select 6 vertices, and among these there are ${6\choose 3}\cdot \tfrac12$ ways to choose a $K_{3,3}$, since one must choose $3$ vertices of the 6 and each $K_{3,3}$ is counted exactly twice (choosing vertices $\{1,2,3\}$ or $\{4,5,6\}$ result in the same graph). \\
		
		\vspace{50px}
		
		\item How many subgraphs of $K_{10}$ are isomorphic to $C_{10}$?\\
		
		\noindent There are $10!$ ways to order the 10 vertices, but each $C_{10}$ is counted 20 times: once for each starting point of the cycle (of which there are 10 possibilities) and twice for each direction: for example, the cycle $(1,2,3,4,5,6,7,8,9,10)$ can be represented
		$$
		(1,2,3,4,5,6,7,8,9,10), (1,10,9,8,7,6,5,4,3,2), (2,3,4,5,6,7,8,9,10,1), (2,1,10,9,8,7,6,5,4,3), \dots 
		$$
		so the answer is $10!/(20) = 9!/2$.
		
		\vspace{50px}
		
		\item Show that if $K_{10}$ is expressed as an ear decomposition, the number of ears (including the cycle at the beginning) is the same regardless of the decomposition, and compute this number. \\
		
		The initial cycle has the same number of edges and vertices. Each ear after that introduces $k$ new edges and $k-1$ new vertices, for some $k\geq 1$. Thus, because $K_{10}$ has $10$ vertices and $45$ edges, there must be exactly 35 ears after the initial cycle, hence $36$ ears total.
	\end{enumerate}
	
	
	\newpage \textbf{Question 3}: Recall that an \textit{Eulerian cycle} on $G$ is a sequence of vertices $v_1,v_2,\dots,v_n$ such that $v_i\sim v_{i+1}$ for each $i$, every edge of $G$ is $(v_i,v_{i+1})$ for exactly one $i$, and $v_1=v_n$.
	
	\begin{enumerate}[(a)]
		\item Show that $K_4$ has no Eulerian cycle. If you use a fact proven in class or the textbook, you must prove it. \\
		
		The vertices of $K_4$ all have degree 3. If an Eulerian cycle enters $v$ via edge $a$ and leaves via edge $b$, then $v$ cannot be visited again because there is only one remaining edge; once it is used, there is no way to leave $v$ without traversing a previously used edge. Thus there is no Eulerian cycle.
		\vspace{60px}
		
		\item Show that $K_5$ \textit{does} have an Eulerian cycle by constructing an explicit example. \\
		
		Let the vertices of $K_5$ be labeled $1,2,3,4,5$. One Eulerian cycle is
		$$
		(1,2,3,4,5,1,3,5,2,4,1).
		$$
		\vspace{60px}
		
		\item Show \textit{using induction} that $K_n$ has an Eulerian cycle for every odd $n\geq 3$.\\
		
		The base case, $n=3$, is fairly easy. So assume that $K_n$ has an Eulerian cycle, where $n$ is odd, and let us show it for $K_{n+2}$. Let the vertices of $K_{n+2}$ be $v_1,v_2,\dots,v_n,a,b$, and let $E$ be an Eulerian cycle on $v_1,\dots,v_n$ (by the inductive hypothesis), starting and ending on $v_1$. Then the following is an Eulerian cycle of $K_{n+2}$:
		$$
		(\underbrace{v_1, \dots, v_1}_{E}, a, v_2, b, v_3, a, v_4, b, \dots , v_n, a, b, v_1).
		$$
		Note that $v_n$ is followed by $a$ \textit{because $n$ is odd}.
		
		\vspace{50px}
		
		\item Is every graph with an Eulerian cycle 2-connected? Prove it or give a counterexample.\\
		
		No. A counterexample is two triangles which share a vertex $v$. If this vertex is removed, then the graph is disconnected, so it is not 2-connected. But it does have an Eulerian cycle, formed by joining the cycles of the two triangles which start and end at $v$. To be specific, if the triangles are $v,a_1,a_2$ and $v,b_1,b_2$, then we have the cycle $v,a_1,a_2,v,b_1,b_2,v)$.
	\end{enumerate}
\end{document}